\documentclass[11pt]{article}
\newcommand{\LabNumber}{10}
\newcommand{\LabTitle}{Maps, Hash Tables, and Binary Search Trees}
\newcommand{\LabTopic}{Map ADT, hash functions, collision handling, BST ordered map}
\input{common.tex}

\begin{document}
\labheader

\section*{Objectives}
\begin{itemize}[leftmargin=*,nosep]
  \item Define the Map ADT and implement a simple list-based map for reference.
  \item Design a hash function and implement separate-chaining collision resolution.
  \item Implement \texttt{ChainHashMap<K,V>} with $O(1)$ expected \texttt{get}/\texttt{put}/\texttt{remove}.
  \item Implement a \texttt{BSTMap<K,V>} ordered map with $O(h)$ operations and inorder iteration.
\end{itemize}

\begin{summary}[Quick Reference]
\textbf{Map interface.}
\begin{lstlisting}
V get(K key);           // null if absent
V put(K key, V value);  // returns old value (or null)
V remove(K key);
int size(); boolean isEmpty();
Iterable<K> keySet(); Iterable<V> values();
\end{lstlisting}

\textbf{Hash function (polynomial rolling).}
\begin{lstlisting}
int hashCode(String s) {
    int h = 0;
    for (char c : s.toCharArray()) h = 37 * h + c;
    return h;
}
int compress(int hash) { return Math.abs(hash) % capacity; }
\end{lstlisting}

\textbf{Separate chaining.}\quad Each bucket is a \texttt{LinkedList<Entry<K,V>>}; scan with \texttt{key.equals()}.\\
\textbf{Load factor} $\lambda = n/N$: rehash (double) when $\lambda > 0.75$.\\
\textbf{BST ordered map.}\quad \texttt{put/get/remove}: $O(h)$; $h = O(\log n)$ average, $O(n)$ worst.
\end{summary}

\subsection*{Quick Experiments (Try It)}

\begin{tryit}{Hash collisions}
Hash \texttt{"Aa"} and \texttt{"BB"} with Java's \texttt{hashCode()}.
Do they collide? Try other two-character pairs.
What does this tell you about compression and bucket chaining?
\end{tryit}

\begin{tryit}{Load factor and rehashing}
Insert entries one by one into your \texttt{ChainHashMap}.
Print the load factor after each insert and note when rehashing occurs.
\end{tryit}

\section*{In-Lab Exercises (target: 2 hours)}

\begin{labexercise}{List-Based Map \hfill {\small \itshape (20 min)}}
Implement \texttt{ListMap<K,V>} backed by a singly linked list of \texttt{Entry<K,V>}.
\begin{itemize}[leftmargin=*,nosep]
  \item \texttt{put}: scan for existing key, update or prepend.
  \item \texttt{get}/\texttt{remove}: linear scan with \texttt{key.equals()}.
\end{itemize}
All operations $O(n)$. This is your $O(n)$ baseline.
\end{labexercise}

\begin{labexercise}{\texttt{ChainHashMap<K,V>} \hfill {\small \itshape (45 min)}}
Implement using separate chaining, initial capacity 11 (prime).
\begin{itemize}[leftmargin=*,nosep]
  \item Compression: \texttt{Math.abs(key.hashCode()) \% capacity}.
  \item \texttt{put}/\texttt{get}/\texttt{remove}: find bucket, scan with \texttt{equals}.
  \item Rehash when load factor $> 0.75$ (double to next prime).
  \item \texttt{keySet()}: iterate all buckets.
\end{itemize}
Test: insert 20 entries, look up 5, remove 3, print remaining keys.
\end{labexercise}

\begin{labexercise}{\texttt{BSTMap<K extends Comparable<K>, V>} \hfill {\small \itshape (40 min)}}
Implement an ordered map backed by a BST:
\begin{itemize}[leftmargin=*,nosep]
  \item Recursive \texttt{put}, \texttt{get}, \texttt{remove} (three-case delete with inorder successor).
  \item \texttt{K firstKey()} and \texttt{K lastKey()} (leftmost/rightmost).
  \item \texttt{Iterable<K> keySet()}: inorder traversal yields sorted keys.
\end{itemize}
Test with the same 20 entries. Verify \texttt{keySet()} output is sorted.
\end{labexercise}

\begin{labexercise}{Word Frequency Counter \hfill {\small \itshape (15 min)}}
Count word frequencies in a 50-word paragraph (case-insensitive, strip punctuation with
\texttt{replaceAll("[{\textasciicircum}a-zA-Z]","")}).
Use \texttt{ChainHashMap}, then \texttt{BSTMap}. Print the top-5 most frequent words.
Verify both maps give the same counts.
\end{labexercise}

\begin{aicollab}
\textbf{Suggested prompts:}
\begin{itemize}[leftmargin=*,nosep]
  \item ``Compare \texttt{HashMap} and \texttt{TreeMap} in Java: when would you choose each?''
  \item ``Why is a prime number recommended for hash table capacity?''
\end{itemize}
\medskip\textbf{Verify:} Ask the AI to predict bucket distribution for keys 0--9 in a table of
capacity 7. Check with your implementation.
\end{aicollab}

\takehomebanner

\begin{trace}[Hash table trace]
Insert keys 10, 22, 31, 4, 15, 28, 17 into a hash table of capacity 7 ($k \bmod 7$).
Draw the table (separate chaining). Which buckets have chains of length $> 1$?
\end{trace}

\begin{trace}[BST trace]
Insert: 20, 10, 30, 5, 15, 25, 35. Draw the BST.
Trace \texttt{get(15)}: list each node visited.
Trace \texttt{remove(10)} (two-child case): show the inorder successor replacement.
\end{trace}

\begin{writecode}[Open addressing: linear probing]
Write \texttt{ProbeHashMap<K,V>} using linear probing.
On collision at index $i$, try $i+1, i+2, \ldots$ (mod capacity).
Use a \emph{defunct sentinel} for removes so future probes are not broken.
Explain how \texttt{get} skips defunct entries but stops at a null slot.
\end{writecode}

\vspace{4pt}
\noindent\textbf{Submission.} Save files as \texttt{lab10\_ex*.java} and submit on the course site.

\end{document}
